% Part: first-order-logic
% Chapter: syntax-and-semantics
% Section: introduction

\documentclass[../../../include/open-logic-section]{subfiles}

\begin{document}

\olfileid{fol}{syn}{int}

\olsection{引言}

为了建立一阶逻辑的理论与元理论，我们必须首先定义其表达式的语法和语义。
一阶逻辑的表达式分为项和公式。
项由变元、常项符号和函数符号形成。
公式则由谓词符号与项一起形成（它们形成最小的、“原子”公式），然后我们可以利用逻辑联结词和量词，从原子公式出发构造更复杂的公式。
设定形成规则的方式有很多种；我们只给出其中一种可能的方案。
其他系统会选择不同的符号，选取不同的联结词集合作为原始联结词，采用不同的括号使用方式（甚至完全不用括号，例如所谓的波兰记法）。
然而，所有方法的共同之处在于：形成规则是以\emph{归纳}的方式定义项和公式的集合。
如果设定得当，每个表达式根据形成规则本质上只能以唯一一种方式生成。
由归纳定义产生的表达式是\emph{唯一可读的}，这意味着我们可以使用同样的方法——归纳定义——来赋予这些表达式意义。

赋予表达式意义属于语义学的范畴。
语义学的核心概念是在结构中的\emph{满足}关系。
结构为语言的基本构件赋予意义：论域是一个非空的对象集合。
量词被解释为以这个论域为取值范围，常项符号被指派以论域中的元素，函数符号被指派以从论域到其自身的函数，谓词符号被指派以论域上的关系。
论域连同对基本词汇的指派一起构成了一个\emph{结构}。
变元可能出现在公式中，为了给出语义，我们还必须给它们指派论域中的元素——这称为\emph{变元指派}。
最终，\emph{满足关系}将这些要素结合起来。
一个公式可以在结构~$\Struct{M}$ 中相对于变元指派~$s$ 被满足，记作 $\Sat{M}{!A}[s]$。
这一关系也是通过对~$!A$ 的结构进行归纳来定义的，例如，利用逻辑联结词的真值表，根据 $!A$ 和~$!B$ 是否被满足来定义 $!A \land !B$ 的满足条件。
由此可知，如果公式~$!A$ 是一个语句（即没有自由变元），那么变元指派就是无关紧要的，因此我们可以直接谈论语句在结构中被满足（或不被满足）。

基于语句的满足关系 $\Sat{M}{!A}$，我们可以进而定义有效性、语义后承和可满足性这些基本的语义概念。
一个语句是有效的（$\Entails !A$），如果每个结构都满足它。
一个语句集语义地后承一个语句（$\Gamma \Entails !A$），如果每个满足~$\Gamma$ 中所有语句的结构也都满足~$!A$。
一个语句集是可满足的，如果存在某个结构同时满足其中的所有语句。
由于公式是归纳定义的，而满足关系又是根据公式的结构归纳定义的，我们可以使用归纳法来证明语义所具有的性质，并建立所定义的各个语义概念之间的关联。




\end{document}